\chapter{Page Layout}
\label{ch:page-layout}

OmniLaTeX uses KOMA-Script's native page geometry system rather than the
\pkg{geometry} package.  This ensures full compatibility with KOMA's
\verb|\areaset|, running headers, and margin-par layout.

\section{KOMA Page Geometry}

KOMA-Script controls page layout through two class options:

\begin{description}
  \item[{\texttt{DIV}}] -- Divisor that determines the text area proportions.
        Higher values yield larger margins.  Common choices: \texttt{DIV=9}
        (narrow margins), \texttt{DIV=12} (default), \texttt{DIV=15}
        (generous margins).
  \item[{\texttt{BCOR}}] -- Binding correction (Bindekorrektur).  Compensates
        for the inner margin lost to binding.  E.g.\ \texttt{BCOR=5mm} shifts
        the text block 5\,mm toward the outer edge on the spine side.
\end{description}

\begin{verbatim}
\documentclass[DIV=12, BCOR=5mm]{omnilatex}
\end{verbatim}

OmniLaTeX defaults to \texttt{DIV=12, BCOR=0mm}.  The \verb|\areaset|
command (KOMA-native) is used internally because it correctly handles the
interaction between text width, text height, and the type area within
KOMA's header/footer system.  Do \emph{not} load the \pkg{geometry}
package alongside OmniLaTeX--it will conflict with KOMA's layout
calculations.

\section{Custom Margins}

For situations that require precise control over all four margins,
OmniLaTeX provides:

\begin{verbatim}
\setCustomMargins{25mm}{20mm}{30mm}{25mm}
%                  left  right  top   bottom
\end{verbatim}

This command converts explicit margin values into equivalent KOMA
\texttt{DIV}/\texttt{BCOR} parameters so that the layout remains
consistent with the rest of the KOMA machinery.  Call it in the preamble
after \verb|\begin{document}| or in a file loaded via
\verb|\BeforeBeginDocument|.

\section{Two-Side Layout}

The \texttt{twoside} class option enables double-sided typesetting:

\begin{verbatim}
\documentclass[twoside]{omnilatex}
\end{verbatim}

In two-side mode:
\begin{itemize}
  \item Odd (recto) pages have a wider right margin; even (verso) pages
        have a wider left margin (mirrored).
  \item Running headers and footers differ between odd and even pages.
  \item \texttt{BCOR} adds extra space on the spine side of every page.
\end{itemize}

One-sided layout (the default) gives uniform margins on all pages.

\section{Headers \& Footers}

OmniLaTeX uses \pkg{scrlayer-scrpage} for running headers and footers:

\begin{verbatim}
\usepackage[automark]{scrlayer-scrpage}
\end{verbatim}

The \texttt{automark} option automatically places the current chapter and
section in the header.  By default:
\begin{itemize}
  \item \textbf{Odd pages:} chapter title (right) and section title
        (left).
  \item \textbf{Even pages:} section title (left) and chapter title
        (right).
  \item \textbf{Page number:} centred in the footer.
\end{itemize}

A thin rule below the header is enabled with \verb|\headerrule|.  Customise
the rule thickness:

\begin{verbatim}
\setheadrule{0.4pt}
\end{verbatim}

\section{Landscape Pages}

For wide tables or figures that do not fit in portrait orientation,
OmniLaTeX provides a \texttt{landscape} environment:

\begin{verbatim}
\begin{landscape}
  \begin{table}
    \caption{Wide data table}
    % ... wide tabular ...
  \end{table}
\end{landscape}
\end{verbatim}

The environment switches to landscape orientation, re-centres the header
and footer, and increments page numbers normally.  For a single-page
landscape section, \verb|\pdflandscape| can be used as a command.

\section{A5 Format}

OmniLaTeX supports A5 output via the \texttt{a5} class option:

\begin{verbatim}
\documentclass[a5paper]{omnilatex}
\end{verbatim}

All margins, font sizes, and float placements scale appropriately.  A5
is useful for pocket-sized handouts or supplementary booklets.

\section{Page Numbering}

LaTeX and KOMA-Script provide three standard numbering schemes:

\begin{description}
  \item[{\texttt{\textbackslash frontmatter}}] -- Roman numerals (i, ii, iii, \ldots).
        Used for the title page, abstract, table of contents, and preface.
  \item[{\texttt{\textbackslash mainmatter}}] -- Arabic numerals (1, 2, 3, \ldots).
        Used for the body of the document.  Resets to 1.
  \item[{\texttt{\textbackslash backmatter}}] -- Roman numerals again.
        Used for the bibliography, appendices, and index.
\end{description}

\begin{verbatim}
\frontmatter
\tableofcontents
\mainmatter
\chapter{Introduction}
\backmatter
\bibliography{refs}
\end{verbatim}

\section{Intentionally Blank Pages}

In two-side mode, chapters open on recto (odd) pages.  KOMA-Script
inserts a blank verso page when needed.  OmniLaTeX marks these pages with
the \texttt{blankpage} page style, which suppresses headers and footers
but prints a small note at the bottom:

\begin{verbatim}
\renewcommand*{\blankpage}{%
  \thispagestyle{empty}%
  \null\vfill
  \centering\itshape This page is intentionally left blank.
  \vfill}
\end{verbatim}

To disable the note, set \verb|\blankpagestyle{empty}| in the preamble.
